\documentclass[11pt,a4paper]{article}
\usepackage[margin=22mm]{geometry}
\usepackage{amsmath,amssymb,amsthm,booktabs,microtype,needspace}
\usepackage[T1]{fontenc}
\usepackage{lmodern}
\usepackage[colorlinks=true,linkcolor=blue,urlcolor=blue]{hyperref}
\hypersetup{pdftitle={Subset sums: a sharper non-progression lemma and the audited bound 94},pdfauthor={Sum-of-Subset research project}}
\newcommand{\FS}{\operatorname{FS}}
\newcommand{\T}{\mathcal T}
\newcommand{\Z}{\mathbb Z}
\newtheorem{theorem}{Theorem}[section]
\newtheorem{lemma}[theorem]{Lemma}
\newtheorem{proposition}[theorem]{Proposition}
\title{Subset Sums in a Fourfold Sumset\\\large A sharper non-progression lemma and the audited bound $N\le94$}
\author{Sum-of-Subset research project}
\date{16 October 2026}
\begin{document}
\maketitle
\begin{abstract}
We continue the independently audited chain from $N\le95$ to $N\le94$.
The main methodological improvement precedes the endpoint computation:
we prove $\rho(X)\ge |X|-5$ for a positive real set with non-progression
subset-sum support and $|X|\ge17$. This strengthens the previously used
$|X|-14$ bound and lowers the archived sharp lemma's threshold from 21
to 17. The finite base has 2,723 deleted intervals, each independently
checked by two exact support computations. At the size-95 endpoint it
eliminates the entire positive-low-hole three-empty family and removes
10,280 low-hole root systems before geometric solving.
Loss-dependent inequalities are selected by explicit group-wide
threshold tests. The complete endpoint proof uses fresh integer
certificates; no archived endpoint certificate is read.
The stronger archived bound 84 remains outside the complete independent
replay: this report advances the audited chain and does not surpass 84.
\end{abstract}

\section{Statement, notation, and dependencies}
For a finite real set $A$, let
$\FS(A)=\{\sum_{x\in X}x:X\subseteq A\}$ include the empty sum and let
$4A=A+A+A+A$ allow repeated summands. Let $N$ be the largest feasible
cardinality for $\FS(A)\subseteq4A$.
\begin{theorem}\label{main}
Every feasible set has at most 94 elements. Thus $71\le N\le94$.
\end{theorem}
The lower bound is the unchanged set
$[-31,31]\cap\Z\cup\{\pm93,\pm155,\pm217,\pm713\}$.
We use the audited chain $115\to103\to100\to99\to98\to97\to96\to95$,
ending in \cite{previous}. Its integralization, covering, peeling,
small-increment inverse, simultaneous inverse, span, effective-hole,
sharp four-extreme, mass, and pairing lemmas remain the foundation.
The new defect argument develops the archived strategy \cite{archive},
with a larger independently enumerated base and a smaller threshold.

A hypothetical size-95 example contains zero, since adjoining zero
otherwise produces a forbidden size-96 example. No hereditary
feasibility under arbitrary deletion is assumed. Integralize, remove
four original outer elements on each side to obtain $D$, and peel twice
more to $D\supset D_1\supset E=P\cup\{0\}\cup(-Q)$. Put
\begin{gather*}
 d=87,\quad e=79,\quad m=a+b=78,\quad a=|P|,\quad b=|Q|,\\
 T_r=r(r+1)/2,\quad I=(a-39)^2,\quad
 \Lambda=1+T_a+T_b=1561+I,\quad\Omega=450-I.
\end{gather*}
For a positive set $X$ of size $r$, write
\[
 \mu(X)=|\FS(X)|-1-T_r,\quad
 \delta_0(X)=|\{x\in X:2x\notin R(X\cup\{0\})\}|,\quad
 \rho(X)=\mu(X)+2-2\delta_0(X),
\]
where $R$ denotes distinct-element two-sums. The inherited bounds from \cite{foundation} are
$\rho\ge-8$ in general, $\rho\ge0$ for $r\ge14$, and
$\mu\ge r-3$, $\rho\ge r-14$ for a non-progression support in the
ranges used below. For increasing prefixes $X_i$, the increments
$y_i=|\FS(X_i)|-|\FS(X_{i-1})|-i$ are nonnegative integers, and
\begin{equation}\label{recurrence}
 \rho(X_i)\ge\rho(X_{i-1})+y_i-2.
\end{equation}
Indeed, $\mu$ increases by $y_i$ and the number of diagonal defects
increases by at most one. If $i\ge6$ and $y_i\le i-5$, the inherited
small-increment inverse expresses $X_i$, up to positive scaling, as
$[1,i+q]_{\Z}\setminus H$, where $q\le y_i$ and the maximum is retained.

\section{A sharp defect bound starting at size 17}
\begin{theorem}\label{sharpnonap}
Let $X$ be a positive real set of size $r\ge17$. If $\FS(X)$ is not an
arithmetic progression, then $\rho(X)\ge r-5$. The constant 5 is sharp.
\end{theorem}
We first prove the deleted-interval input, including its complete finite
base. The computer checks in this section form an explicit finite part
of the proof; the following reduction supplies the unbounded part.

\begin{lemma}\label{deleted}
If $Y=[1,j+q]_{\Z}\setminus H$, $j\ge8$, $0\le q\le3$, $|H|=q$,
and $j+q\notin H$, then a non-progression $\FS(Y)$ satisfies
$\rho(Y)\ge j-5$.
\end{lemma}
\begin{proof}
The density forces $\gcd(Y)=1$, so an arithmetic-progression support is
exactly a full interval. The case $q=0$ is full. For $j\ge10$,
$j+q\le3j/2-2$, and Lev's central-interval theorem
\cite[Theorem 1]{Lev1998} gives
\[
 [2q+1,\sum Y-(2q+1)]_{\Z}\subseteq\FS(Y).
\]
By subset-complement symmetry, the total number of holes is twice the
number of missing sums in $[1,2q]$. These small sums depend only on
which positions through $2q$ are deleted.

Every diagonal defect lies in $[1,q]$ or $[j,j+q]$. For an interior
$x$, the $\min(x,j+q-x)>q$ distinct pairs centered at $x$ in
$[0,j+q]$ are disjoint, so $q$ deletions cannot destroy them all.
Low defects depend only on deleted positions through $2q$; high defects
depend only on deleted depths through $2q$ below the maximum.

For $j\ge14$, preserve these low positions and high depths, with
respective counts $l,u$. Replace the other $t=q-l-u$ deleted depths
by $2q+1,\ldots,2q+t$ in an interval with 14 retained elements. These
positions fit between the two endpoint bands for every $q\le3$.
The resulting set $Y_0$ has the same number of support holes and the
same number of diagonal defects. With $h=j+q$, the exact identity is
\[
 \mu(Y)=\sum_{v\in H}(h-v)-T_{q-1}-\operatorname{holes}(\FS(Y)).
\]
Each low deletion contributes $j-14$ more to this depth sum than in
$Y_0$. High depths do not change, and the ordered middle depths are at
least the replacement depths. Non-full support requires $l\ge1$:
otherwise every integer through $2q$ is present and there are no holes.
Consequently
\[
 \rho(Y)-j\ge\rho(Y_0)-14+(l-1)(j-14)\ge\rho(Y_0)-14.
\]
It remains to check exactly $j=8,\ldots,14$, $q=1,2,3$, and every
$q$-subset of $[1,j+q-1]$. The new checker
\path{verification/nonap_sharp.py} computes subset sums both as a set
of integers and as an integer bitset, compares the results, and forms
all distinct-pair sums of $Y\cup\{0\}$ to count defects. The counts are
\begin{center}\small
\begin{tabular}{lrrrrrrr}\toprule
$j$&8&9&10&11&12&13&14\\\midrule
All deletion sets&164&219&285&363&454&559&679\\
Non-full supports&88&108&130&154&180&208&238\\
Minimum $\rho(Y)-j$&$-5$&$-5$&$-5$&$-5$&$-5$&$-5$&$-5$\\\bottomrule
\end{tabular}
\end{center}
All 2,723 cases, including all 1,106 non-full supports, pass. The checker
also verifies the depth identity and defect bands in every base case.
It checks the central-interval hole identity only for $j\ge10$;
that theorem is not assumed at sizes 8 or 9. The preceding signature
reduction now proves the assertion for every $j\ge8$.
\end{proof}

\begin{proof}[Proof of Theorem~\ref{sharpnonap}]
If every $y_i\ge4$ for $8\le i\le r$, start at
$\rho(X_7)\ge-8$ and use \eqref{recurrence}:
\[
 \rho(X_r)\ge-8+2(r-7)=2r-22\ge r-5\qquad(r\ge17).
\]
Otherwise take the last $j\ge8$ with $y_j\le3$. Since $3\le j-5$,
the small-increment inverse applies. If $\FS(X_j)$ is not a progression,
Lemma~\ref{deleted} gives $\rho(X_j)\ge j-5$; all later increments
increase $\rho$ by at least two, proving the result.
If the support is a progression, let $k>j$ be the first index whose
support is not. The preceding progression and its new translate must
be disjoint: an intersection places the translating element on the same
step and makes their union another progression. Hence
\[
 y_k=|\FS(X_{k-1})|-k\ge1+T_{k-1}-k,
 \qquad\rho(X_k)\ge T_{k-1}-k-9\ge k-5
\]
for $k\ge9$. Apply the same later-increment argument.
For sharpness, take $X=\{1,3,4,\ldots,r+1\}$. Lev's theorem with
$q=1$ shows that the only missing sums in its hull are 2 and
$\sum X-2$. Its only diagonal defects are 1 and $r+1$; for 3 use
$0+6$, and for other nonmaximal elements use their two neighbors.
Thus $\mu=r-3$ and $\rho=r-5$.
\end{proof}
Compared with the archived threshold-21 proof, adding the 383 cases at
sizes 8 and 9 permits the recurrence to start two steps earlier and
lowers the threshold by four. We make no claim that 17 is the optimal
threshold; it is the threshold proved and used here.

\section{Lattice reduction without a circular use of the new budget}
The nonnegative inherited identity is
\begin{equation}\label{budget}
 \rho(P)+\rho(Q)+2\vartheta+c+\T=\Omega,
 \quad c=2|2D|+18d+85-|\FS(D)|,
 \quad\T=\kappa+\ell_1+\ell_2+g,
\end{equation}
where $\kappa=|\FS(E)|-\Lambda-\mu(P)-\mu(Q)$ and the other terms
are the same nonnegative peeling and diagonal slacks as in \cite{previous}.
A deficient sign would give $1+T_{79}>24\cdot95+77$. The paired
bootstrap gives $q^2\le2416$, with even $q$, hence $q\le48$ and
$a,b\ge15$. Nonnegative $\Omega$ then leaves $a=18,\ldots,60$.
For all 43 oriented pairs, $T_{\min(a,b)}-2-\Omega\ge68>0$.
The simultaneous inverse theorem therefore supplies one common core
lattice, with holes
\[
 h_P,h_Q\le\kappa,\quad z=h_P+h_Q-\kappa,\quad
 0\le z\le\min(h_P,h_Q)\le\kappa,\quad h_X<|\FS(X)|-1.
\]
This step has not assumed that $D_1$ or $D$ is on that lattice.

The inherited universal span comparison, with
$C=2m+4e+14-\Lambda=-1075-I$, gives
\begin{align*}
 |\FS(D_1)|-2|\FS(E)|&\le C+\Omega+\rho-\mu+2z,\\
 |\FS(D_1)|-2|\FS(E)|&\le C+3\Omega-\rho-\mu,
\end{align*}
where $\rho,\mu$ denote the two side sums. The maximum of a positive
set is defective, so $\rho\le\mu$. If $z\le14$, the first bound is
at most $-597-2I<0$. If $z\ge15$, both supports are non-progressions.
Even the earlier bounds $\rho\ge50$, $\mu\ge72$ leave only
$|a-39|\le6$. These sides have size at least 33. The effective-hole
witness $(r,C_0,\beta)=(33,89,7)$, checked with strict integer tests,
gives $\mu\ge78+180=258$. Thus $\rho+\mu\ge308>275-4I$.
The maximum resulting high-hole comparison is $-33$.
The first-layer comparison has margin
$\Lambda-(4d-10+\Omega)=773+2I>0$.
An off-lattice element would give two disjoint translates of the smaller
subset-sum support, contradicting the strict factor-two comparisons.
Thus $D_1$ and then $D$ are on the core lattice.

Only after this reduction do hole propagation and the sharp
four-extreme theorem of \cite{previous} give
\begin{equation}\label{holebudget}
 w\le4z,\qquad \T\ge\kappa+164-3z,
 \qquad \T\ge\max(55,z,164-2z).
\end{equation}
The new defect theorem is not needed for these lattice comparisons;
its use below is therefore independent of the sharper post-lattice
budget. All core sides now have at least 18 elements, so its
threshold 17 is always satisfied.

\section{Scalar pruning and common endpoint models}
Let $j$ of the eighteen covering units be empty, and put
$\delta=c-87j$. The universal span estimate combined with
\eqref{budget}--\eqref{holebudget} yields
\[
 R_p+R_q+\delta\le368-I-87j+4z.
\]
For zero holes, $c\le286-I$. For $1\le z\le14$, both sides are
non-progressions and Theorem~\ref{sharpnonap} gives
$\rho\ge(a-5)+(b-5)=68$. Therefore
\[
 c\le218-I+2z\le246-I<3d.
\]
In particular there are no positive-low-hole three-empty modes.
The complete low-hole ranges and counts are
\begin{center}\small
\begin{tabular}{lrrrrrrr}\toprule
Regime&$Z_0$&$Z_1$&$Z_2$&$Z_3$&$P_0$&$P_1$&$P_2$\\\midrule
$a$&23--55&25--53&29--49&34--44&24--54&27--51&31--47\\
Integer roots&132&2030&12138&32912&124&1750&9826\\\bottomrule
\end{tabular}
\end{center}
Their total is 58,912. The weaker $\rho\ge50$ would retain 69,192
low roots, including 8,976 positive three-empty roots. The sharper
lemma removes those roots and 1,304 additional positive roots,
a total reduction of 10,280 before any linear feasibility computation.
The zero-hole three-empty family is retained in full.

For high holes the effective-hole search is the same bounded search as
in \cite{previous}, with every chosen parameter rechecked exactly.
It gives lower bounds $u_p,u_q$ for $\mu(P),\mu(Q)$. Replace the defect
floor in the flexible inverse by the sharper floor:
\[
 r_*(r,u)=\max\{r-5,u+2-2r,L_2(r,u),\ldots,L_{10}(r,u)\},
\]
where, for $v=\max(0,u-4\beta+2)=(\beta+1)q+s$,
\[
 L_\beta(r,u)=\min\{(\beta-1)(r-14),(\beta-1)q+\min(s,\beta-1)\}.
\]
The constant $r-14$ inside this inherited flexible-inverse formula is
not silently changed. Thus
\[
 c\le450-I-r_*(a,u_p)-r_*(b,u_q)-\max(55,z,164-2z).
\]
There are 3,500 surviving oriented scalar states and $c\le171<174$,
so at most one unit is empty. The active-center estimate uses the
orientation fixed by the first pair's actual contributor, exactly as
in \cite{previous}; selecting a stronger reflected orientation is not
allowed. Against $F\le1692+4z$ it leaves 959 states for $j=0$.
There are 1,371 states with enough cap for $j=1$.

For each fixed $(a,j,u_p,u_q)$, take the largest hole count $Z$ and
largest cap $C_*$. Replace the two mass coordinates by
$\widehat S_p=S_p+Z-z$, $\widehat S_q=S_q+Z-z$.
All mass, cardinality, span, loss and grouped-overlap rows retain their
validity by the transport lemma of \cite{previous}; upper endpoint rows
weaken. The 2,330 retained states merge into 50 groups: 15 of type
$H_0$ and 35 of type $H_1$. This reduces high-hole integer models from
$959\cdot4+1371\cdot70=99806$ to $15\cdot4+35\cdot70=2510$.
For low positive holes the same transport uses $Z=14$.

\section{Threshold-driven rows and the unified covering budget}
The twenty centers, eighteen paired units, actual-contributor choices,
and the grouped union-deficit variables are exactly those proved in
\cite{previous}. For $j=0,1,2,3$, their mode counts are
$4,70,578,2992$. All missing-unit subsets and both choices in a potentially
split pair are retained. Let $l_v,u_v$ be the hull losses at a retained
center, and let $L_p,L_q$ be the actual-target losses of the two leading
slots. The common residual budget contains these losses and the
nonnegative actual deficits of disjoint groups of slot labels.

For clarity, for $k$ ordinary same-sign slots the core-hull lower bound
on their joint deficit is
\[
 79k-(R_p+R_q+v_{\max}-v_{\min}+1)+\max_v(l_v+u_v).
\]
For a leading slot and $k$ own-center slots the corresponding lower
bound is
\[
 79(k+2)-2-p_{\max}-R_p-2R_q+max(L_p,l_1+u_1,\ldots,l_k+u_k).
\]
Both also dominate the sum of individual losses. These follow by
counting the common core hull and then the noncore points, and adding
points of the raw union outside the actual target. Each whole-block
variable is the actual deficit of that block, so alternative subset
bounds constrain one variable; they are not separate overlap gains
added repeatedly. Sum only deficits of disjoint blocks.
The leading centers are integral with at most one empty unit by the
inherited contributor lemma. With two or three empty units,
$\delta\le112<173\le|2D|$, so both leading slots contribute and their
centers are again integral. This justifies all core-hull counts.

\subsection{Every conditional row has its own checked threshold}
For each structural group $g$ the checker derives its maximum residual
cap $C_g$ from its entire integer side-size interval, or from its high
cap $C_*$. It also takes the smallest side size in that group.
The following tests are independent:
\begin{center}\small
\begin{tabular}{ll}\toprule
Rows used&Required group-wide condition\\\midrule
Ordinary individual quantiles&$C_g<\min(a,b)+5$ throughout the group\\
Core individual quantiles&$C_g<83$\\
Basic leading quantiles&$C_g<173$\\
Nested leading quantiles&$C_g<165$\\
Strict outer order&$j\le1$, or $j=2$ with contributing leading slots\\\bottomrule
\end{tabular}
\end{center}
Use the strongest justified individual quantile type; otherwise use
global quantiles. This avoids two invalid substitutions from the
preceding endpoint: an $H_1$ cap can now be 84, so core quantiles are
not automatic; an $H_0$ cap can be 171, below 173 but above 165, so basic
leading rows do not automatically justify nested ones. All three-empty
models retain weak outer order. $Z_3$ has residual cap at most 25 and
uses ordinary quantiles; $P_2$ has cap at most 72 and uses core quantiles.
The code computes these choices, rather than accepting a certificate's
choice of applicability conditions.

During high-hole transport, freeze the weaker endpoint group's row
policy when comparing a source state with its endpoint. A smaller
source cap could justify additional rows, but those rows are not needed
for the common endpoint model. This keeps the matrix and row labels
fixed and proves coefficientwise transport dominance without conflating
the two leading thresholds.

\subsection{Numerical constants in the necessary models}
Write $K=|2D|$, $x_0=R_p<x_1<x_2<x_3<x_4=M_p$, and similarly $y_i$.
The inherited mass and gap rows on each side are
\begin{align*}
 R_p&\ge a,&M_p-R_p&\ge4,&x_i-x_{i-1}&\ge1,\\
 2M_p+6R_p-2S_p+2z&\le-2T_a-12-2u_p,\\
 M_p+4R_p-S_p&\le-T_{a-1}-6,\\
 x_1+x_2+x_3+M_p-S_p+z&\le-T_a-u_p,\\
 R_p+x_1+x_2+x_3+M_p-S_p&\le-T_{a-1}.
\end{align*}
Here $u_p=0$ for $Z$, $a-3$ for $P$, and the certified high-hole
lower bound for $H$; all occurrences of $z$ are transported together.
The common cardinality row is
\[
 S_p+S_q+\delta-2K-4z\le1650-87j.
\]
At a positive center $v$ the global quantiles are
$v+l_v+M_p\ge87$, $v-u_v-M_q\le S_p-86$.
The core versions replace $M_p,M_q$ by $R_p,R_q$ and 87,86 by 83,82;
the ordinary versions are $v+l_v\ge b+5$, $v-u_v\le S_p-a-4$.
Individual loss caps are 86 globally and 82 for core rows. Reflection
supplies the negative side.

The basic leading rows are
\[
 c_p+2M_p+L_p\ge K,\quad S_p+2M_q-c_p+L_p\ge K-1,
 \quad S_p+L_p\ge K.
\]
Their coarse versions use $K\ge173$. The nested rows, used only under
their stronger cap test, are
\[
 c_p+2x_i+L_p\ge165+2i,\quad
 S_p+2y_i-c_p+L_p\ge164+2i\qquad(0\le i\le4).
\]
They count the contributors from a subset of $D$ of size $e+4+i$,
whose twofold sumset has at least $2(e+4+i)-1=165+2i$ elements.
The cap below 165 ensures the counting interval is nonempty.

The coarse model substitutes
$K\le M_p+M_q+R_p+R_q+21$ into the cardinality row, eliminating $K$
and the inner coordinates. Difficult groups retain all five inherited
nested interval-union envelopes, plus the 25 independent-cutoff rows
\[
 K\le M_p+M_q+x_i+y_h+1+T_{4-i}+T_{4-h}\quad(0\le i,h\le4).
\]
The latter count the integer central interval and all pairs among the
remaining same-sign outliers. No new unproved twofold-sumset estimate
is inserted.

\subsection{Exhaustive local alternatives}
The four types of binary alternatives from \cite{previous} remain valid,
with current constants. For a retained center, either $l_v+u_v\ge83$,
or $l_v+u_v\le82$ and both core quantiles hold. For ordinary endpoints,
either $v+l_v\ge b+5$, or $l_v\ge b+5$ and $v\le-1$; either
$v-u_v\le S_p-a-4$, or $u_v\ge a+5$ and $v\ge S_p+1$.
For each inner rank $i=1,\ldots,4$, either $l_v\ge i$ or
$v\ge y_{5-i}+1$, and either $u_v\ge i$ or $v\le S_p-x_{5-i}$.
Finally a potentially split pair has either the necessary condition
for its other contributor, or its selected actual member is sole and
its two hull losses can be added to the individual-loss budget, with
global quantiles 87 and 86.

These are alternatives for actual configurations, whose losses and
centers are integral. The model relaxes the coordinates to nonnegative
reals. After transport an upper-endpoint ordinary alternative is
selected using the enlarged endpoint; the original upper loss still
dominates the lost-point count. A necessary positivity condition is
never used as a sufficient condition for contribution. Every binary
node retains both children; no hypothesis limiting the number of
large-loss centers is imposed.

\section{Complete exact coverage and reproducibility}
The verifier reconstructs the scalar domain, the complete covering-mode
grammar, and every row policy before reading the root assignments.
It checks the finite base of Theorem~\ref{sharpnonap} on every run.
A certificate cannot omit the new lemma's finite cases or enable a
conditional row merely by naming it. The final endpoint counts are
\begin{center}
\begin{tabular}{lr}\toprule
Quantity&Count\\\midrule
Structural root groups&6,806\\
Integer root systems&61,422\\
Coarse / nested structural groups&6,758 / 48\\
Direct symbolic interval checks&6,809\\
Extremum-point evaluations&16,611\\
Integer roots requiring trees&182\\
Internal binary nodes / tree leaves&493 / 675\\
Maximum tree depth&10\\
Final integer leaf systems&61,915\\
Distinct integer certificate vectors&897\\\bottomrule
\end{tabular}
\end{center}
The split counts are 342 core-or-large, 109 ordinary-or-outside,
11 inner-endpoint, and 31 both-or-sole pair splits. The identities
$675=182+493$ and $61915=61422+493$ are checked.

For $Ax\le b$, $x\ge0$, the certificate supplies nonnegative integer
row weights $w$ with $w^TA\ge0$ and $w^Tb<0$. Therefore
$0\le w^TAx\le w^Tb<0$, a strict contradiction. The minimum normalized
gap $-w^Tb/\sum_iw_i$ is $1/91$. For each side-size interval, the matrix
is fixed and the right-hand sides are directly constructed quadratic
polynomials. Exact arithmetic checks their weighted maximum at the
endpoints and the integers around the vertex. Every parameter is assigned
exactly once; missing roots, unknown rows, unproved or repeated splits,
and incomplete trees are rejected.

The 19 negative controls include the invalid one-large-loss shortcut,
universal core rows for the $Z_2$ or $H_1$ families, nested leading rows
at cap 171, a stronger false $r-4$ defect bound, use of the new theorem
below its proved threshold, arbitrary reflection in the active estimate,
and an unsupported exclusion of cardinality 94. This verifier proves
$N\le94$, not $N\le93$.

\paragraph{Verification files.}
The following commands need only Python 3.10+ and its standard library.
The endpoint data are \path{data/upper94_certificates.json}:
\begin{verbatim}
python -I -S -B verification/nonap_sharp.py
python -I -S -B verification/upper94.py --output .work/upper94.json
python -I -S -B verification/upper94.py --exhaustive
\end{verbatim}
Ordinary mode checks symbolic intervals and all tree leaves. Exhaustive
mode additionally rebuilds every integer model, compares every row with
its symbolic expression, and checks all assigned weights pointwise.
The finite lemma checker recomputes all 2,723 cases and their inequalities;
its enumeration digest is an audit record, not a substitute for checking.

The optional SciPy search in
\path{experiments/generate_upper94_certificates.py} saves only candidates
that pass exact checking. Results, the exceptional-group index, logs and
receipts are in \path{results/2026-10-16/upper94/}. Standard and optimized
Python outputs are byte-identical; exhaustive and interval results agree
apart from the mode label; regenerated certificates are byte-identical;
and upper-95 regression is unchanged. The Windows entry
\path{verify_upper94_windows.bat} runs all eight finite continuations.

A program PASS checks the finite proof data. The unbounded
deleted-interval reduction and real-set argument are proved above;
they and the cited foundation have not been formalized in a proof assistant.

\paragraph{Next endpoint.}
At size 94, $d=86$, $e=78$, $m=77$, and the balanced parameters
$\Omega=467$, $G=162$ allow a positive low-hole cap
$467-67-162+28=266\ge3\cdot86$. The positive three-empty family
reappears. That endpoint needs new lattice, threshold and coverage checks.

\begin{thebibliography}{9}
\raggedright
\bibitem{previous} Sum-of-Subset research project,
\emph{Joint covering losses and the audited bound $N\le95$},
\path{reports/upper95_from96_20261016.tex}, 16 October 2026.
\bibitem{foundation} Sum-of-Subset research project,
\emph{The audited continuation from 115 to 103},
\path{reports/upper103_from115_20261014.tex}, 14 October 2026.
\bibitem{archive} Sum-of-Subset research project,
\emph{Complete upper-86 research report},
\path{reports/upper86_complete_20261005.tex}, 5 October 2026.
\bibitem{Lev1998} V.~F.~Lev, \emph{On consecutive subset sums},
Discrete Mathematics \textbf{187} (1998), 151--160, Theorem 1.
\url{https://doi.org/10.1016/S0012-365X(98)80006-X}.
\end{thebibliography}
\end{document}
